\chapter{Introduction}
\label{ch:introduction}

\section{Background and Motivation}

Colleges continuously invest in preparation and assessment materials:
syllabi, textbooks, notes, practice questions, and examinations. In most
institutions these assets are created and managed with generic office tools and
file storage. The consequences are well known: teaching material is duplicated
or lost, syllabus meaning has to be manually re-derived for every subject,
question banks accumulate without validation of correctness or alignment with
exam patterns, and assessments are distributed and evaluated manually, so
students receive little or no structured feedback.

The project team built a platform that addresses these problems as \textbf{a
multi-tenant, role- and permission-controlled web application} for educational
institutes. A single deployment serves many institutes on one shared codebase,
yet keeps every institute's data isolated and inaccessible to the others. The
platform makes the preparation-and-assessment workflow structured: uploaded
materials are OCR-processed into reusable text, syllabus documents are analysed
into confirmed curriculum structures, study content is generated
semi-automatically, question banks are built against explicit paper patterns,
and assessments are scheduled, taken, and evaluated server-side.

\section{Problem Statement}

\noindent The problem addressed by this project is stated as follows:

\emph{Educational institutes require a single, tenant-isolated platform that (i)
converts unstructured teaching material (scanned PDFs, images, documents) into
structured, searchable content through OCR and AI-assisted processing; (ii)
turns syllabus documents into confirmed academic structures; (iii) builds
curated question banks and examination assets from declarative paper patterns;
(iv) runs scheduled assessments with automatic server-side evaluation and
meaningful analytics; (v) offers ungraded practice for students; and (vi)
exports all assets in standard, print-ready formats --- all while restricting
every operation to the acting user's role, permissions, and academic scope.}

\section{Objectives}

\noindent The specific objectives of the implementation are:

\begin{enumerate}
  \item \textbf{Tenancy and identity:} model institutes, memberships, roles, and
  users, and isolate data per tenant.
  \item \textbf{Academic structure:} model years, classes, divisions, subjects,
  class--subject offerings, and teacher/student assignments, and derive each
  actor's academic scope from that structure.
  \item \textbf{Materials and OCR:} upload PDF, image, and document materials
  with chunked upload, process them through a distributed OCR pipeline with
  embedded-text-first extraction and PaddleOCR fallback, and expose per-page
  corrections.
  \item \textbf{Syllabus intelligence:} analyse syllabus documents into a
  confirmed global curriculum with topics, units, objectives, and learning
  outcomes.
  \item \textbf{Aided content creation:} generate notes, flashcards, Cornell
  notes, summaries, and important-concepts content from eligible source
  materials through an async AI generation pipeline, with provenance.
  \item \textbf{Question assets:} establish a data-driven question-type system
  and a question bank with generation, approval states, ranking, and difficulty
  distribution; build paper patterns and question papers from blueprints.
  \item \textbf{Examination:} author assessments from blueprints, publish and
  activate them on schedules, record attempts with immutable question
  snapshots, evaluate objective answers server-side, and produce aggregate
  analytics.
  \item \textbf{Practice:} provide flashcard and question practice modes with
  instant per-item feedback and no grading.
  \item \textbf{Export:} generate print-ready PDF, DOCX, and XLSX outputs for
  content, questions, papers, and results.
  \item \textbf{Security and architecture:} enforce tenant isolation at the
  data layer, security decisions from an explicit permission catalogue, and
  strictly internal compute for OCR and AI via a provider gateway.
\end{enumerate}

\section{Scope}

The implemented system covers the full lifecycle described above. It does not
implement online-answer-sheet (FORM), OMR-scanning, or on-screen-marking
checking features, automated grading of subjective (essay-type) answers, or
per-attempt N-of-M mechanics; these are identified as deferred future work in
Chapter~\ref{ch:conclusion}.

\section{Organisation of the Report}

\noindent Chapter~\ref{ch:literature-review} reviews related systems and
technologies. Chapter~\ref{ch:requirements} presents the requirements
analysis. Chapter~\ref{ch:methodology} describes the development methodology and
the project timeline. Chapter~\ref{ch:architecture} describes the system
architecture and Chapter~\ref{ch:database-design} the database design.
Chapter~\ref{ch:security} details the security and authorization architecture,
and Chapter~\ref{ch:features} explains the core features and their workflows.
Chapter~\ref{ch:implementation} covers implementation details,
Chapter~\ref{ch:testing} the testing and validation strategy, and
Chapter~\ref{ch:results} the results and discussion.
Chapter~\ref{ch:conclusion} concludes the report and lists future scope.